\section{The round-cone exponent}
\label{sec:rt}

We retain the round-cone experiments, their regularity tests, and their
free growth exponents from the preceding section.  We write
$\Gamma_{\mathrm{cl},\mathrm{round}}$ and
$\Gamma_{\mathrm{wav},\mathrm{round}}$ for the classical and wave
exponents, respectively.  The allowance in the following theorem is
deliberately larger than the losses in its proof.

\begin{theorem}[Round-cone bound]\label{thm:round}
For every $0<\eta<1$, the round-cone experiments satisfy
\[
 \Gamma_{\mathrm{cl},\mathrm{round}}\le 100\eta,
 \qquad
 \Gamma_{\mathrm{wav},\mathrm{round}}\le 100\eta.
\]
The same upper exponent bounds hold with a fixed finite list of admissible
row contractions, with arbitrary separable Hilbert multiplicity, and with
the finite-power support and error conventions of the round-cone model.
\end{theorem}

The assertion concerning contractions follows from
Proposition~\ref{rt:free-composition}.  We prove the free
assertions, first classically and then for waves.  In a contradiction
argument, $\Gamma$ denotes a saturated exponent with
\[
 \Gamma>100\eta,
 \qquad
 \Gamma>\Gamma_{\mathrm{cl},\mathrm{round}}
 \quad\hbox{in the wave case}.
\]
As in the preceding section, take the least homogeneous mass exponent
$p$ in the free closure at $\Gamma$, and put
\begin{equation}\label{eq:rt-pd}
 \Gamma=d+\frac{1-p}{2},\qquad
 p\ge1,\qquad 0<d\le1.
\end{equation}
At a prescribed depth $a$, the energy and regularity benchmarks are
\begin{equation}\label{eq:rt-benchmarks}
 E_a=e s^{-da+o(1)},\qquad
 \kappa_a=\kappa s^{pa+o(1)}.
\end{equation}
Throughout this section a finite list of depths is fixed before the
inner scale limit.  An assertion with $s^{o(1)}$ means its strict-power
version with arbitrarily small prescribed error, after the earlier
extremal and homogeneous approximations have been taken sufficiently
accurate.  We never construct an infinite stack of operators.

The thin-shape reduction leaves round labels. In the wave case we place
them at the uncertainty scale. We then exclude concentration of velocities
near a plane and use compatible names and stationarity to obtain an
outgoing entropy bound. Direct integration closes the classical argument;
for waves, the bound supplies the two differential inequalities used in
the final descent.

\subsection{Round labels occur at the uncertainty scale}

Proposition~\ref{rt:thin-exclusion}, together with
Theorem~\ref{thm:cylinder}, shows that every homogeneous ending family
which saturates the assumed exponent has $f=s^{o(1)}$.  This statement
also applies to a raw peeling family read before synthesis: insert a
short positive free gap immediately before that readout and apply
Equation~\eqref{eq:CE} on the gap.

\begin{lemma}[Volume of a classical round test]
\label{lem:rt-test-volume}
Normalize the parent length to one and let $\nu$ be a classical
phase measure supported in a fixed bounded box of $(b,y,A)$, with
Lebesgue density at most $L_0$.  Uniformly in the centers and
coefficients of every admissible round or hyperquadric test,
\begin{equation}\label{eq:rt-test-volume}
 \frac{\nu(\text{test})}{r^4f^{1-\eta}}
 \le C L_0\min(r^2,r^{-3})f^\eta.
\end{equation}
The constant depends only on the support box and the fixed test
coefficient bound.
\end{lemma}

\begin{proof}
Use the normalized test coordinates $(x,z,V)$ at its initial time.
The phase Jacobian is $r^6$: horizontal position contributes $r^2$,
vertical position contributes $r^2$, and velocity contributes $r^2$.
For a hyperquadric test, $F$ is independent of $V$, and
\[
 g=\partial_VG=\alpha+2\beta V-2kx,
 \qquad |g|^2=1+4\beta G-4kF.
\]
The coefficient bound is $100$.  On the two sublevels, if
$f\le1/1600$, the identity gives $|g|^2\ge1/2$.
At least one of its two components then has magnitude at least
$1/2$.  Partition the angular sublevel according to such a
component.  Hold the other angular coordinate fixed.  The function
$G$ in the remaining coordinate is a quadratic, so the portions
where its derivative has magnitude at least $1/2$ consist of at
most two monotonic intervals.  On each interval the preimage of
$[-f,f]$ has length at most $4f$.  Integrating over the bounded
other coordinate gives normalized angular area $O(f)$, uniformly
at each fixed $(x,z)$ satisfying $|F|\le f$.  Integration over
$|x|,|z|\le1$ bounds the physical phase volume by $Cr^6f$.

For a second bound, hold the physical position $(b,y)$ fixed and
integrate only over the fixed bounded support box in $A$.  On the
same sublevel, $\partial_A G=g/r$, so the identical quadratic
slicing argument gives angular area at most $Crf$; the other
physical angular coordinate ranges in a fixed bounded interval.
Integration over the compact physical position support yields
phase volume at most $Crf$.  These arguments do not restrict the
test centers: only the points satisfying its normalized round
conditions and lying in the physical support are being counted.

For $f\ge1/1600$, the crude phase volume is at most
$C\min(r^6,1)$, which is bounded by
$C\min(r^6,r)f$ after enlarging the constant.  This also handles
round tests, for which $f=1$.  Combining the two volume bounds,
multiplying by $L_0$, and dividing by the weight proves
Equation~\eqref{eq:rt-test-volume}.
\end{proof}

\begin{proposition}[Angular uncertainty at a saturated wave cut]
\label{prop:rt-planck}
In the wave contradiction, a homogeneous saturating family at any
positive-depth cut $j$ satisfies
\begin{equation}
 \frac{r_j}{d_{\ell_j}}=s^{o(1)}.
 \tag{round-Pl}\label{eq:round-Pl}
\end{equation}
After normalizing the terminal angular width to $r_1=1$, it follows that
\begin{equation}
 H=s^{1+o(1)},\qquad u(a)=\frac{1-a}{2},\qquad
 r_a=s^{u(a)+o(1)}.
 \tag{round-u}\label{eq:round-u}
\end{equation}
\end{proposition}

\begin{proof}
The uncertainty condition and $f=s^{o(1)}$ give the lower bound in
Equation~\eqref{eq:round-Pl}.  Suppose the ratio is large by a fixed
power.  Choose an entering free gap of ratio $\Delta=s^{h+o(1)}$, with
$h>0$ smaller than that power and smaller than the available gap.
Normalize the parent time and position lengths to one and the ending
angular width to one.  By decreasing $h$, we can arrange
$d_\Delta\le\Delta^M$ for any fixed $M$.  In these units write
$\kappa_*$ for the parent regularity.

Partition the entering plateau rows into bounded phase boxes, each with
a bounded angular range after a boost.  Equation~\eqref{eq:acc} shows
that an ending label receives contributions from at most a subpower
number of these boxes.  Coherent splitting therefore has only a
subpower cost.  The raw mass $e_g$ of one box satisfies
$e_g\le s^{-o(1)}\kappa_*$ by round regularity, and its synthesized
field retains the same regularity upper bound.

Regard the raw row measure in this box as a positive measure of exact
rays.  Average it over independent smooth position translations and
velocity boosts, of size $\Delta^2$, using the accompanying exact
vertical tilt.  More explicitly, at the initial time the transformation
with parameters $(u,v,a)\in\R^2\times\R\times\R^2$ is
\[
 (b,y,A)\longmapsto(b+u,\ y+v+2a\cdot(b+u),\ A+a).
\]
At subsequent times the velocity boost additionally changes $b$ by
$ta$ and $y$ by $t|a|^2$, so it sends exact rays to exact rays.
For each original phase point, the displayed map from $(u,v,a)$ to
the new phase point has Jacobian one.  Smooth parameter densities
at scale $\Delta^2$ therefore give a jittered density at most
$C\Delta^{-10}e_g$, with bounded support.
These symmetries preserve admissible test bounds exactly, not just
up to an enlargement.  Transform the test center and its central
velocity by the same symmetry.  Substitution in the normalized
coordinates leaves $(x,z,V)$ unchanged, including their time
extension.  Thus $r,f$, the normalized polynomial coefficients,
their discriminant, and their bound of $100$ are unchanged.
The uncertainty parameter $d_1$ is also unchanged.  Pullback of a
test with $rf\ge d_1$ is consequently an admissible input test
of the same weight; averaging its bound preserves regularity.
For a test below that uncertainty threshold,
Lemma~\ref{lem:rt-test-volume} gives
\[
 \frac{\mu(\text{test})}{r^4f^{1-\eta}}
 \le C\Delta^{-C}e_g\min(r^2,r^{-3})f^\eta.
\]
If $rf<d_1\le\Delta^M$, the final factor is at most
$\Delta^{M\eta}$: for $r\le\Delta^M$ use $r^2$; for
$\Delta^M<r\le1$ use $f\le\Delta^M/r$; and for $r\ge1$ use
$r^{-3}f^\eta\le\Delta^{M\eta}r^{-3-\eta}$.
Choosing $M>C/\eta$ therefore gives admissible classical regularity
$s^{-o(1)}\kappa_*$.  This choice depends on $\eta$ and the fixed
jitter construction, not on the original polynomial clipping powers.

At each ending time, group the wave assignments into spatial cubes of
side $\Delta$.  Equation~\eqref{eq:cone-kernel}, applied to the
complement of the ray preimage of an enlarged cube, and contraction on
that preimage imply that its analyzed wave mass is bounded by the raw
positive ray mass of the enlarged cube, up to negligible error.
Here the packet position errors are much smaller than $\Delta$.
Jitter moves a ray by $O(\Delta^2)$ over the bounded time interval, so
all its jittered copies belong to one further enlargement.  These
enlarged cubes have bounded overlap, have bounded angular width, and
are round tests of bounded weight.  Splitting a label among cubes and
combining assignments in the same cube cost only $s^{-o(1)}$.

Consequently the classical round-cone bound controls the ending wave
cube sum on this gap, relative to $e_g\sqrt{\kappa_*}$, with exponent
$\Gamma_{\mathrm{cl},\mathrm{round}}+o(1)$.  The kernel errors are
first made smaller than all required powers on the clipped inner
sequence; components of negligible energy are estimated directly on
the polynomial family of output tests.  Summing the raw component
budgets contradicts saturation of a gap whose exponent is strictly
larger.  This proves Equation~\eqref{eq:round-Pl}.  Substituting
$d_{\ell}=\sqrt{H/\ell}$, first at the terminal cut and then at depth
$a$, proves Equation~\eqref{eq:round-u}.
\end{proof}

\subsection{A short step with velocities near a plane}

We establish the estimate that will remove planar concentration from
conditional velocity laws.  Normalize a free parent interval to length
one.  The selected angular spectrum is bounded after a boost; ending
round labels have angular width one and time mesh
$\Delta=s^{h'+o(1)}$.  Their weights are $s^{o(1)}$ and will be omitted
in this subsection.  Let $e_g$ and $\kappa_*$ be the input energy
budget and regularity.  In the wave case assume
$H=\Delta s^{o(1)}$.  A plane slab in velocity space has the form
\begin{equation}
 |d_0+\alpha\cdot A+\beta|A|^2|\lesssim\rho^\eps,
 \qquad |\alpha|^2+\beta^2=1,
 \qquad \rho=s^h,
 \tag{plane}\label{eq:plane}
\end{equation}
where $h,\eps>0$ are fixed.  If the slab meets the bounded angular
range, then $|d_0|\lesssim1$.  Classical input rays are supported in
this slab.  For waves the spatial spectrum is supported within
$s^{-o(1)}\sqrt\Delta$ in velocity of its bounded portion, in the
prescribed positive compact radial bands.

\begin{lemma}[Spatial readout]\label{lem:rt-spatial-readout}
For a free wave with this bounded angular spectrum,
\begin{equation}
 \Delta^{-1/2}\sum_\lambda m_\lambda^{3/2}
 \lesssim s^{-o(1)}\int_{\mathcal U}\norm{u(t,b,y)}^3\,dt\,db\,dy
       +\text{negligible error}.
 \tag{space-3}\label{eq:space-3}
\end{equation}
Here $\mathcal U$ can contain the bounded time interval and all space.
If the contributing input rows are localized in bounded phase boxes,
only their subpower enlarged spacetime region is needed.
\end{lemma}

\begin{proof}
Group the ending assignments at each cut $T$ in ordinary spatial
$\Delta$-cubes.  A label meets only $s^{-o(1)}$ such cubes.  The
analyzer at scale $\Delta$ and propagation for $|t-T|\le\Delta$ have
spatial kernel tails at scale $\Delta s^{-o(1)}$: rescale their compact
smooth symbols, using $H=\Delta s^{o(1)}$ and bounded angular
support, and integrate by parts.  Thus a cube mass $m_D$ is at most
the spatial square energy in an enlarged cube $D^+$ at each
$t\in[T,T+\Delta]$, plus negligible error.  On $D^+$, Holder's
inequality gives
\[
 m_D^{3/2}
 \lesssim s^{-o(1)}\Delta^{3/2}
               \int_{D^+}\norm{u(t,b,y)}^3\,db\,dy.
\]
Average over this time interval, multiply by $\Delta^{-1/2}$, and
sum over the bounded-overlap spacetime cubes.  The same argument with
polynomial clipping and kernel truncation proves the localized
statement.  All errors can be made negligible relative to an original
reference energy before estimating components of very small energy.
\end{proof}

\begin{proposition}[Plane-slab estimates]\label{prop:rt-plane}
There is a universal constant $C$ with the following properties.
Let $\lambda>0$ be sufficiently small relative to $\eps$.
If
$|\alpha|^2-4\beta d_0\lesssim\rho^{2\lambda}$, one may choose
$0<h'<\lambda h$, inside any prescribed available positive gap, so
that
\begin{equation}
 \Delta^{-1/2}\sum_\lambda m_\lambda^{3/2}
 \le s^{-o(1)}\sqrt{\kappa_*}\,e_g.
 \tag{tiny-cap}\label{eq:tiny-cap}
\end{equation}
In the complementary case choose a fixed
$0<h'\le\eps h/20$ inside the available gap, then take
$\lambda$ sufficiently small relative to $h'/h$.  With
$P=\rho^{-\lambda}$,
\begin{equation}
 \Delta^{-1/2}\sum_\lambda m_\lambda^{3/2}
 \lesssim s^{-o(1)}P^C\Delta^{-4\eta}
                    \sqrt{\kappa_*}\,e_g.
 \tag{sec}\label{eq:sec}
\end{equation}
The estimates hold classically even after further positive
restrictions or raywise contractions.  For waves they concern the free
continuation of the specified input state.
\end{proposition}

\begin{proof}
Put $D_0=|\alpha|^2-4\beta d_0$.  At an approximate zero of the
polynomial in Equation~\eqref{eq:plane}, its squared angular gradient
is $D_0+O(\rho^\eps)$.  If $D_0\lesssim\rho^{2\lambda}$ and
the slab is nonempty, $|\beta|$ is bounded below: otherwise
$|\alpha|$ is bounded below and $D_0$ is bounded below.  Completing
the square shows that all slab velocities lie in a ball of radius
$O(\rho^\lambda)$.  Our choice $h'<\lambda h$ makes that radius
$O(\Delta)$.

Classically, after boosting at its center, the preimage of an ending
spatial $\Delta$-cube has angular and horizontal widths
$O(\Delta)$ and tilted vertical width $O(\Delta)$.
It is covered by $O(\Delta^{-1})$ round phase boxes of radius
$O(\Delta)$, of total weight $O(\Delta^3)$.  Hence
$\max m_\lambda\le s^{-o(1)}\kappa_*\Delta^3$.
Since the total assigned mass over all times is at most
$s^{-o(1)}\Delta^{-1}e_g$, Equation~\eqref{eq:tiny-cap} follows
from $\sum m^{3/2}\le(\max m)^{1/2}\sum m$.

In the wave case the boosted horizontal and vertical frequency
widths are at most $s^{-o(1)}\Delta^{-1/2}$ and
$s^{-o(1)}\Delta^{-1}$.  A spatial box of dimensions
$(\sqrt\Delta,\sqrt\Delta,\Delta)$ at a bounded time can receive
energy only from the correspondingly enlarged parent phase box, up
to kernel tails.  Round regularity at radius $\sqrt\Delta$ and
contraction bound this energy by
$s^{-o(1)}\kappa_*\Delta^2$.  Reproduce the field with a smooth
band kernel.  Cauchy--Schwarz against that kernel, partitioning its
tails into enlargements of these boxes, gives
\[
 \norm{u(t,b,y)}^2\lesssim s^{-o(1)}\kappa_*.
\]
An arbitrarily small fixed-power enlargement followed by arbitrary
kernel decay handles the tails on each inner sequence.  Combine this
bound with energy conservation and
Lemma~\ref{lem:rt-spatial-readout} to obtain
Equation~\eqref{eq:tiny-cap}.

We turn to $D_0\gtrsim\rho^{2\lambda}=P^{-2}$.  The exact
section has angular gradient at least a constant multiple of $P^{-1}$,
and each approximate solution is within $O(P^C\rho^\eps)$ of it.
Cover it by $P^C$ buffered patches.  Boost an exact direction in each
patch to zero, rotate, and divide the plane equation by its nonzero
transverse derivative.  The section then reads
$A_2+\widetilde\beta|A|^2=0$, with
$|\widetilde\beta|\lesssim P$.  Apply the linear change
\begin{equation}
 b'_2=b_2+\widetilde\beta y,\qquad
 t'=t+2\widetilde\beta b_2+\widetilde\beta^2y,\qquad
 b'_1=b_1,\qquad y'=y.
 \tag{L}\label{eq:L}
\end{equation}
This preserves $ty-|b|^2$.  On a ray set
\[
 D_A=1+2\widetilde\beta A_2+
                       \widetilde\beta^2|A|^2.
\]
The transformed direction satisfies
\[
 \frac{dt'}{dt}=D_A,\qquad
 A'_1=\frac{A_1}{D_A},\qquad
 A'_2=\frac{A_2+\widetilde\beta|A|^2}{D_A},\qquad
 |A'|^2=\frac{|A|^2}{D_A}.
\]
Choose the patch to have sufficiently small inverse-polynomial radius
so that $D_A$ is bounded above and below.  Its derivatives and those
of the inverse chart have bounds $P^{C(n+1)}$ at order $n$.
Our choices of $h'$ and then $\lambda$ ensure that
$|A'_2|\le\Delta^8$ classically and
$|A'_2|\le P^C\sqrt\Delta$ for waves.  The patches have buffers
larger than the finitely many admitted spectral broadenings.

Partition also into bounded input position charts.  Each output label
knows only $s^{-o(1)}P^C$ of these input charts, by
Equation~\eqref{eq:acc}.  Restrict the parent rows and synthesize in
the wave case.  The sum of raw energy budgets is bounded by the
original budget times this overlap; inherited regularity and spectral
support have the stated bounds.  Thus it suffices to work in one
chart, where all relevant transformed coordinates and times have
size at most $P^C$.

We first record its input regularity.  At a fixed transformed time,
test $(A'_1,b'_1,y')$ in widths $(R,R,R^2)$ with the usual tilt by
twice the base angular center, and test $(A'_2,b'_2)$ in widths
$(W,W)$, with the angular center of $A'_2$ equal to zero.  For
$0<R\lesssim1$ and $W\le P^C R$, the parent mass on the predicted
rays satisfying these conditions is at most
\begin{equation}
 P^C\kappa_* R^{3+\eta}W^{1-\eta}.
 \tag{slice-reg}\label{eq:slice-reg}
\end{equation}
For waves the widths are enlarged, if necessary, so that $R,W$ are
at least the parent uncertainty width.  To verify the bound, pull
back the round box and backflow to the parent time.  The null
covector
$dy-2A\cdot db+|A|^2dt$ is preserved up to a controlled scalar;
Taylor's formula therefore gives horizontal errors $P^C R$ and
tilted vertical errors $P^C R^2$.  The pullback of
$b'_2-b'_{2,*}$ is the polynomial
$b_2+\widetilde\beta y-b'_{2,*}$, of discriminant one.
Its value and ray derivative at the parent time are bounded by
$P^C W$, because their values at the tested time have that bound
and the propagation time is at most $P^C$.  In angular $R$-units,
divide this polynomial by $R$.  Lemma~\ref{rt:covering} then
applies with thickness $P^C W/R$ and gives
$P^C\kappa_*R^4(W/R)^{1-\eta}$.  If $W>R$, the round bound
alone is enough.  Boxes larger than a chart buffer are covered by
$P^C$ boxes, and all coefficient bounds are still polynomial in $P$.
This proves Equation~\eqref{eq:slice-reg}.

For the classical estimate, put $W=\Delta^8$ and partition at
$t'=0$ into $W$-slabs of $b'_2$.  In one slab, project onto the base
rays with velocity $(X,X^2)$, $X=A'_1$.  Omitting $(A'_2)^2$
from the vertical velocity changes positions by at most
$P^C W^2$, negligible at all scales $R^2$ with
$R\gtrsim\Delta^2$.  Equation~\eqref{eq:slice-reg} implies the
hypothesis of Equation~\eqref{eq:kin} with constant
$P^C\kappa_*W^{1-\eta}$, since $R^\eta\le1$.
The affine change of spacetime compares old and new
$\Delta$-cells with $P^C$ multiplicities.  A ray meeting a new
spacetime cell meets an enlarged spatial cell at its nearest grid
time.  Apply Equation~\eqref{eq:kin} on the $P^C$ required unit time
intervals and use Cauchy--Schwarz with total counted mass at most
$P^C\Delta^{-1}$ times the slab mass.  The sum of $3/2$ powers in
this slab is at most
$s^{-o(1)}P^C\sqrt{\kappa_*W^{1-\eta}}$ times its mass.
A full spatial cell meets at most $P^C\Delta/W$ slabs; convexity
costs the square root of this number.  After multiplication by
$\Delta^{-1/2}$ the remaining factor is
$W^{-\eta/2}=\Delta^{-4\eta}$.  This proves
Equation~\eqref{eq:sec} classically.

For waves the change in Equation~\eqref{eq:L} produces an exact
transformed wave $u'$.  The polar covector transforms with
$\rho'=\rho D_A$.  On the buffered patch this is a smooth
invertible change of the spatial cone coordinates, with Jacobian
and inverse bounded by powers of $P$ and $\rho'H$ in a fixed
positive band.  Plancherel gives transformed-time energy at most
$P^Ce_g$.  The same bound holds when only a subset of parent rows
is synthesized: insert a common smooth spectral projection on the
larger patch before making the change of variables.

Choose $W=P^{C_0}\sqrt\Delta$ so that $|A'_2|\le W$.
Cover the required transformed time range by intervals of length
$L=P^{-C_2}$, where $C_2$ is large compared with the fixed bandwidth
powers.  At fixed $b'_2$, multiply by a band-limited Schwartz time
cutoff $\chi((t'-t'_0)/L)$.  The resulting function of
$(t',b'_1,y')$ lies within thickness $O(L^{-1})$ of the base cone:
the additional temporal frequency has size
$\rho'|A'_2|^2\le P^C$.  After scaling lengths by $L$, apply
Equation~\eqref{eq:GWZ} at frequency $L/H$.  The fine angular
windows may depend on $\xi'_1/(2\rho')$ alone, and hence commute
with restriction in $b'_2$ and the time cutoff.  Comparable adapted
windows give the same bound by convolution with rapidly decaying
dual-sector kernels and a summable covering by translated tiles.
At angular width $R$, where
$(H/L)^{1/2}\lesssim R\lesssim1$, the tile volume is
$|U|\asymp L^3R^3$.

The fine-cap square energy on such a tile, at fixed $b'_2$, obeys
\begin{equation}
 \int_U\sum_{\vartheta\subset O(\tau)}
                 \norm{u'_{\vartheta}}^2
 \le s^{-o(1)}P^C\Delta^{-\eta/2}\kappa_*|U|.
 \tag{tile}\label{eq:tile}
\end{equation}
Indeed the second spatial bandwidth is
$B_{\mathrm{freq}}=P^C/\sqrt\Delta$.  Reproduction and
Cauchy--Schwarz in $b'_2$ reduce the left side to
$C B_{\mathrm{freq}}$ times the corresponding energy integrated
over an enlarged $W$-slab.  At each remaining tile time,
Equation~\eqref{eq:slice-reg} bounds the energy of all the relevant
fine caps by $P^C\kappa_*R^{3+\eta}W^{1-\eta}$.
They have a common predecessor set of parent rows; after restriction
to it, synthesis, the spectral change, and Plancherel with bounded
overlap give this bound without a fine-cap counting factor.
Integration over the tile time length $O(L)$ gives, after division
by $|U|$,
\[
 P^C\kappa_*\Delta^{-1/2}W^{1-\eta}R^\eta
 \le P^C\kappa_*\Delta^{-\eta/2}.
\]
Powers of $L^{-1}$ have been included in $P^C$.

For completeness, the predecessor restriction used here follows
directly from the transformed packet kernel.  If $p$ is a parent
row, its starting point becomes $(t'_p,b'_p,y'_p)$, and the phase is
\[
 (b'-b'_p)\cdot\xi'-(y'-y'_p)\rho'
                 -(t'-t'_p)|\xi'|^2/(4\rho').
\]
The duration is at most $P^C$, and angular support lies within
$P^C\sqrt H$ of its transformed direction $A'_p$.  In variables
$H\rho'$ and $\sqrt H(\xi'-2\rho'A'_p)$ the amplitudes and
their derivatives have bounds $P^{C(n+1)}$.  The fine-cap cutoff
obeys the same bounds because its width is at least
$(H/L)^{1/2}\ge\sqrt H$.  Integration by parts restricts the
predicted positions to errors $P^C\sqrt H$ horizontally and
$P^C H$ after tilt.  These lie inside the enlarged conditions of
Equation~\eqref{eq:slice-reg}.  Choose the enlargement power larger
than the derivative power per order; increasing the order then
gives arbitrary decay.  Schur bounds on polynomially clipped rows,
including the polynomial number of caps, turn this into negligible
operator error.  The reproducing kernel in $b'_2$ has arbitrary
decay outside the $W$-slab, because $B_{\mathrm{freq}}W$ can be
made larger than a fixed power of $P$.  The time cutoff supplies
the same decay outside an enlarged time range.  Orders are chosen
after the inner polynomial support and energy ratios are fixed.

Use Equation~\eqref{eq:tile} for one factor in each square on the
right side of Equation~\eqref{eq:GWZ}.  Summing the remaining
factor over tiles and integrating in $b'_2$ costs only the total
spacetime square energy.  Sum the $P^C$ time windows to obtain
\[
 \int_{|t'|\le P^{C_1}}\norm{u'}^4
 \lesssim s^{-o(1)}P^C\Delta^{-\eta/2}\kappa_*e_g.
\]
The arbitrarily small exponent loss in Equation~\eqref{eq:GWZ}
is included in $s^{-o(1)}$.  Cauchy--Schwarz against the spacetime
square energy gives an $L^3$ bound with factor
$P^C\Delta^{-\eta/4}\sqrt{\kappa_*}e_g$.
Changing coordinates back and using
Lemma~\ref{lem:rt-spatial-readout} proves
Equation~\eqref{eq:sec}, with room in its stated loss.

All nonnegligible patch, Jacobian, window, and overlap losses used a
fixed power of $P$.  The kernel orders affect only the negligible
remainders.  Thus the exponent $C$ can be fixed before choosing
$\lambda,h,h'$.  In applications choose the large-discriminant
$h'/h$ first, then $\lambda$, and finally the tiny-cap step.
\end{proof}

\subsection{Finite configurations and compatible names}

Let $\ell_a\sim s^a$ and $R_v\sim s^v$ be dyadically rounded.
The time $T_a$ is the left endpoint of the ancestor interval at depth
$a$.  Denote the lower corner of the dyadic angular $R_v$-cell by
$A_v$.  Define the phase name
\begin{equation}
 C(a,v)=\left(T_a,A_v,
       [b(T_a)]_{\ell_aR_v},
       [y(T_a)-2A_v\cdot b(T_a)]_{\ell_aR_v^2}\right).
 \tag{grid}\label{eq:grid}
\end{equation}
The brackets mean coordinatewise grid cells.  For a wave row measured
at a later cut, the positions in this definition are obtained by
backflow along that row's own ray.  This definition makes no assertion
that a ray is unchanged by synthesis.  The wave domain is
$0<a<1$, $0\le v\le u(a)$; classically we use $v=0$.

\begin{lemma}[Selected laws and rarity transfer]\label{lem:rt-laws}
For any fixed finite rational list of cuts and names, one can choose
the saturated configuration in forward order so that the following
properties hold to arbitrarily small exponent errors.
At every requested subterminal cut $a$, the selected raw mass and
the post-synthesis energy both have size $E_a$.  The raw rows have
intervalwise regularity at most $s^{-o(1)}\kappa_a$ and are carried
on peeling labels with
\begin{equation}\label{eq:rt-weight-count}
 \sum_{\lambda_a}w_{\lambda_a}
       \le s^{-o(1)}E_a/\kappa_a.
\end{equation}
Write $\mathbf P^a$ for their normalized raw row law.  If a portion
tested at a later selected cut $l$ can only be reached through a set
of relative raw mass $q$ at $i<l$, then its relative mass at $l$ is
at most $s^{-o(1)}q^{2/3}$, up to negligible error.  Consequently a
power-rare necessary predecessor remains power-rare at the later
law.  For classical rays the same conclusion holds at the final
assigned law $\mathbf P^1$.
\end{lemma}

\begin{proof}
At each successive cut insert a regularity peel on plateau-analyzed
rows, select an ordinary level which preserves the original terminal
certificate under free continuation, and then synthesize.  Equation~
\eqref{eq:CE}, applied with the fixed finite list, gives
Equation~\eqref{eq:rt-benchmarks} and saturation of the raw and
post-synthesis budgets.  The full tests of the selected level give
Equation~\eqref{eq:rt-weight-count}.  Further preserving
subdivisions need only stay inside these tests; their individual
submasses need not saturate each test.  The choices are made before
later insertions, so the certificate with free continuation remains
available at every chosen cut.

Restrict the $i$-rows to the necessary predecessor set and synthesize.
The resulting states have total energy at most $s^{-o(1)}qE_i$
and regularity at most $s^{-o(1)}\kappa_i$.  The upper cube bound
on the segment gives
\[
 \sum_\lambda m_\lambda^{3/2}w_\lambda^{-1/2}
 \le s^{-o(1)}q\,
       s^{(1/2-\Gamma)(l-i)}E_i\sqrt{\kappa_i}
 =s^{-o(1)}qE_l\sqrt{\kappa_l}.
\]
The identity uses Equation~\eqref{eq:rt-pd}.
Holder's inequality and Equation~\eqref{eq:rt-weight-count} yield
\[
 \sum_\lambda m_\lambda
 \le \left(\sum_\lambda m_\lambda^{3/2}w_\lambda^{-1/2}\right)^{2/3}
      \left(\sum_\lambda w_\lambda\right)^{1/3}
 \le s^{-o(1)}q^{2/3}E_l.
\]
Nonconnections from the complementary rows are negligible by
successive kernel truncations and complementary projections, using
Equation~\eqref{eq:acc}.  In the classical case the restriction is
positive and raywise, so the same proof applies through the final
assignment.  All tested families are polynomial on each inner
sequence; very small components and tail errors are handled before
normalization relative to the original energy.
\end{proof}

We will use deterministic log profiles of finitely many discrete
names.  Their construction is worth specifying.  At a measuring cut
$a$, let $\mathbf P^a_{\rm pre}$ be the normalized raw row law from
Lemma~\ref{lem:rt-laws}, with all earlier choices fixed but before this
additional restriction.  To \emph{preflatten} a fixed finite list of
names and joint names, partition rows by sufficiently fine logarithmic
bins of their probabilities at the sampled values, all computed under
$\mathbf P^a_{\rm pre}$.  Polynomial cardinalities make the
mass with probabilities below a sufficiently small power negligible.
There are only a subpower number of retained logarithmic classes.
The coherent splitting inequality gives one common class $B_a$
preserving the terminal exponent, and Lemma~\ref{lem:rt-laws} forces
$\mathbf P^a_{\rm pre}(B_a)=s^{o(1)}$.  Below, $\mathbf P^a$ denotes
the selected law $\mathbf P^a_{\rm pre}(\,\cdot\mid B_a)$; at any later
preflattening step, its pre-restriction law already includes the earlier
fixed choices.  If a name has mass
$s^{D_*+o(1)}$ before restriction, it has at most
$s^{-D_*-o(1)}$ possible values in this class, and each has mass at
most $s^{D_*-o(1)}$ under the normalized selected law.  Counting
the values whose mass is too small shows that its surprise exponent
equals $D_*+o(1)$ with probability tending to one.  The same
argument for joint names gives conditional scores.  This conclusion
persists after any further restriction of subpower probability,
by likelihood comparison for the numerator and denominator.  At a
fixed strict tolerance the exceptional probabilities can be made
power-small by taking the earlier approximations sufficiently accurate.

\begin{lemma}[Stability of ancestor names]\label{lem:rt-stability}
At a wave cut $a$, preflatten $C(a,v)$ and $(C(a,v),A_w)$, with
$v\le w\le u(a)$.  At a later sampled cut $b$ with
$w\le u(b)$, the same deterministic exponents hold for their
backevaluated names, including after further subpower restrictions
of $\mathbf P^b$.
\end{lemma}

\begin{proof}
On a retained correspondence path, the angular discrepancy between
the two rows is at most $s^{u(b)-o(1)}$.  Backevaluating at $T_a$
gives horizontal error at most $\ell_as^{u(b)-o(1)}$ and vertical
error, tilted by the earlier direction, at most
$\ell_as^{2u(b)-o(1)}$.  These statements follow by adding the
position errors in Equation~\eqref{eq:acc} and the linear and
quadratic errors from changing the velocity during backflow.
The same bounds with the longer time unit apply to earlier anchors.
Since $v,w\le u(b)$, a joint name at either end determines at most
$s^{-o(1)}$ possible joint names at the other end.  When an angular
center changes, recompute its vertical tilt from the horizontal
coordinate: the residual uncertainty is the horizontal grid width
times the angular uncertainty, within the stated vertical width.

The earlier name count therefore bounds the later support, apart
from negligible probability.  A list of too few later names pulls
back to a power-rare list at $a$, because of the individual upper
mass bound there.  Lemma~\ref{lem:rt-laws} excludes a nonnegligible
later probability on such a list.  The count bound excludes
probabilities which are too small.  Apply this to both the root
and joint names to get conditional scores.  Subpower restrictions
preserve these deterministic limits by the preceding likelihood
argument.  On a boundary use slightly coarser widths and then
subpower neighboring lists.
\end{proof}

\subsection{Entropy profiles and time spread}

For a discrete name $U$ and side data $D_*$ write
$i_s(U\mid D_*)=-\log\mathbf P(U\mid D_*)/\log(1/s)$,
evaluated at the random sample.  Side data may be nonatomic; regular
conditional probabilities are used.  Equation~\eqref{eq:LP} and its
finite-list versions will be applied to these scores.

\begin{proposition}[Profiles]\label{prop:rt-profiles}
Passing to finite-list subsequences and then a diagonal, there are
Lipschitz profiles
\begin{equation}
 \begin{split}
 D(a,v)&=\lim\left(i_s(C(a,v))-
                              \log_{1/s}(e/\kappa)\right),\\
 J(a,v,w)&=\lim i_s(A_w\mid C(a,v)),\qquad v\le w\le u(a),
 \end{split}
 \tag{prof}\label{eq:prof}
\end{equation}
in probability on each of the valid measuring laws.  For classical
rays only $D(a,0)$ and $J(a,0,w)$, $0\le w\le1$, are used.
The wave excess satisfies
\begin{equation}
 E(a,v):=D(a,v)-4v-(p+d)a\ge0,
 \qquad E(a,u(a))=0.
 \tag{excess}\label{eq:excess}
\end{equation}
Classically $D(a,0)\ge(p+d)a$ and $D(1,0)\le p+d$.
Furthermore $0\le J(a,v,w)\le2(w-v)$, $J$ is nonincreasing in
$a$ for fixed $v,w$, and in the wave case
\begin{equation}\label{eq:rt-superadditive}
 J(a,v,w)\ge J(a,v,x)+J(a,x,w),\qquad v<x<w.
\end{equation}
Every bounded-branching score increment in these assertions has the
same limit in expectation as in probability.
\end{proposition}

\begin{proof}
A cell $C(a,v)$ restricts one interval to a subpower number of
round tests of radius $s^v$, hence has raw mass at most
$s^{-o(1)}\kappa_as^{4v}$.  Division by $E_a$ gives the lower
bound for $D$ in Equation~\eqref{eq:excess}.  At the wave boundary,
Equation~\eqref{eq:round-u} makes every peeling label determine
$C(a,u(a))$ to a subpower list; its weight is
$s^{4u(a)+o(1)}$.  Equation~\eqref{eq:rt-weight-count} gives the
reverse count bound.  Classically, transfer too-sparse lists to the
final law with Lemma~\ref{lem:rt-laws}; the final weight count gives
the reverse inequality at $a=1$.

Increasing $v$ refines the name up to bounded ambiguity.  An
increment $\delta v$ costs at most $6\delta v$ branching
exponents: two angular, two horizontal, and two vertical, with the
tilt recalculated from the horizontal position.  Increasing $a$
at fixed $v$ likewise refines the name, by backflow within the
angular cell.  It costs at most $4\delta a$ exponents, one for
time and three for position.  The residual tilted vertical velocity
is quadratic in $A-A_v$, which is exactly what is needed for this
backflow comparison.  Adding $A_w$ costs at most two exponents per
unit increase of $w$.  Compare names on a common later measuring
law using Lemma~\ref{lem:rt-stability}; near the domain boundary
change widths first and then the cut.  These comparisons give
uniform Lipschitz bounds for the shifted profiles and their
increments.

Conditional entropy decreases under refinement of the conditioning
name, giving monotonicity in $a$.  Since $C(a,x)$ knows
$C(a,v)$ and $A_x$ up to bounded lists, the entropy chain rule gives
Equation~\eqref{eq:rt-superadditive}.  All entropy comparisons use
bounded-branching increments.  Their tails are uniformly integrable:
if a discrete variable has at most $s^{-B}$ possible refinements,
the probability of surprise exceeding $B+t$ is at most $s^t$.
Thus their expected and typical limiting scores agree, without
requiring a uniform bound on a common unconditional offset.

On each finite rational list choose increasingly accurate
saturation and preflattening data, then take a subsequence on the
bounded shifted scores.  The Lipschitz bounds give a compatible
diagonal and unique continuous extensions.  A finite real query can
subsequently be replaced by rational parameters with arbitrarily
small error in profile values.  This procedure never requires
uniformity in the number of actual gates.  Extra cuts used to test
free continuation from a selected earlier state leave that earlier
state and its established profiles in place.
\end{proof}

\begin{lemma}[Classical angular memory]\label{lem:rt-memory}
For classical rays, at any fixed root name $C(a,0)$ and any
$0<w<1$, the typical conditional surprise of $A_w$ is at least
$(3-p)w-o(1)$.  Consequently
\begin{equation}\label{eq:rt-classical-m}
 m:=3-p=2-2d+2\Gamma\in(0,2].
\end{equation}
\end{lemma}

\begin{proof}
Use the regularized cut $i=1-w$.  Fix a terminal label and an
angular $s^w$-cell.  Their preimage at this cut lies in horizontal
widths $\ell_i s^{w-o(1)}$ and tilted vertical width
$\ell_i s^{w-o(1)}$, including the terminal position error.
It is therefore covered by $s^{-w-o(1)}$ round boxes of radius
$s^w$, of total weight $s^{3w-o(1)}$.
Positive transport and cut regularity bound this mass by
$\kappa_i s^{3w-o(1)}$.  Terminal labels of assigned mass less
than $\kappa_1s^\zeta$ have power-small total mass, for every
fixed $\zeta>0$, by Equation~\eqref{eq:rt-weight-count}.
Dividing by the label mass and using
$\kappa_i/\kappa_1=s^{-pw+o(1)}$ gives the conditional angular
bound given the terminal label.  Such a label determines any
$C(a,0)$ to a subpower list by exact backflow; applying the same
sparse-list bound gives the assertion given $C(a,0)$.
Equation~\eqref{eq:rt-pd} gives the expression for $m$, and the
two-dimensional angular count gives $m\le2$.
\end{proof}

Fix a finite tangent query satisfying
\begin{equation}
 0<a<a+Kh\le b,\qquad K=3,\qquad
 v+Kh<u(b)\quad\hbox{for waves}.
 \tag{query}\label{eq:query}
\end{equation}
Here $b<1$ for waves, whereas $b=1$ and $v=0$ classically.
Set
\[
 C=C(a,v),\quad C'=C(a+Kh,v),\quad
 Z=(A-A_v)/R_v,\quad
 \theta=(T_b-T_a)/\ell_a,\quad \rho=s^h.
\]
Let $W$ be the position at $T_a$, after the central tilt and
subtraction of the bin origins of $C$, divided by its two horizontal
and its vertical grid lengths.  Then $Z,W,\theta$ are bounded,
and the position evolves in these coordinates with velocity
$V(Z)=(Z,|Z|^2)$.  Subscripts below mean ordinary coordinate bins
at width $\rho^z$.  Put $Y_\#=(Z,W)_K$.

\begin{lemma}[Time spread]\label{lem:rt-time}
For every fixed finite list $0<z\le1$,
\begin{equation}\label{eq:rt-time-score}
 i_\rho(\theta_z\mid C,Y_\#)\ge dz-o(1)
\end{equation}
typically.  The statement is power-robust at each strict tolerance,
including after subpower restrictions of the measuring law.
\end{lemma}

\begin{proof}
Use the actual peel cut $l=a+zh$.  Consider a predicted list of at
most $\rho^{-dz+\zeta}$ time cells for each $(C,Y_\#)$, where
$\zeta>0$.  Use Lemma~\ref{cyl:borel-paths} to take a Borel inner
version of the tested list event at $b$, mark the $l$-rows that have
a continuation to it, and retain a Borel inner version of that
analytic mark under the full analyzed $l$-row energy law, inside
the current end-state support. For input restriction,
take the Borel outer hull of all predecessors under the
unrestricted truncated $a$-to-$l$ geometric correspondence, with
the same full analyzed $a$-row energy. A selected row at $a$ reaches at most
$s^{-o(1)}$ possible values of $(C,Y_\#)$ on any such
continuation: concatenate the correspondences and use the strict
uncertainty slack in Equation~\eqref{eq:query}.  Changes of tilt
and bin origins are determined by the available coordinates up to
the same small lists.  A time cell and $T_a$ determine $T_l$ to
bounded ambiguity.  Therefore this earlier row can feed the marked
set at at most $s^{-o(1)}\rho^{-dz+\zeta}$ different $T_l$.
Contraction separately for each such interval bounds total marked
raw mass by
\[
 s^{-o(1)}\rho^{-dz+\zeta}E_a
       =s^{-o(1)}\rho^\zeta E_l.
\]
Complementary nonconnections are negligible.  Transfer the marked
rarity from $l$ to $b$ with Lemma~\ref{lem:rt-laws}.
Testing the popular conditional time cells proves
Equation~\eqref{eq:rt-time-score}.  The same proof with exact
positive paths applies classically.  Earlier errors can be made
smaller than any fixed fraction of $h\zeta$, so the conclusion
persists after a subpower restriction.
\end{proof}

\begin{lemma}[Removal of conditional plane concentration]
\label{lem:rt-slab-removal}
At the query in Equation~\eqref{eq:query}, fix $\eps>0$.
In the wave case one may restrict $\mathbf P^b$ by subpower mass;
in the classical case no additional restriction is necessary.
The resulting laws have all the preceding robust score properties
and satisfy, typically in $(C,C')$,
\begin{equation}\label{eq:rt-plane-nonconcentration}
 \sup_{\Pi}\mathbf P\{\dist(V(Z),\Pi)\le\rho^\eps
                         \mid C,C'\}\longrightarrow0,
\end{equation}
where $\Pi$ ranges over affine planes in $\R^3$.
These choices are made before any conditional product sampling.
\end{lemma}

\begin{proof}
Classically choose one measurable candidate slab for each name pair.
The terminal label determines that pair to subpower lists.
At $i=1-h'>a+Kh$, partition the actual regularized input rays by
the pair and restrict to its specified slab.  Use the two possible
step lengths in Proposition~\ref{prop:rt-plane}, one for each
discriminant range, including both cuts in the finite list.
Choose the large-discriminant $h'/h$ first; then choose $\lambda$
so small that $P^C$ costs less than a fixed part of
$(\Gamma-4\eta)h'$; finally choose the tiny-cap step.
Each case gives a strict gain over exponent $\Gamma$, so for some
$\sigma>0$, possibly different for the two steps,
\[
 \sum m_{\lambda,\mathrm{hit}}^{3/2}
 \le\sum_{h'}s^{-o(1)}
       s^{(1/2-\Gamma+\sigma)h'}E_i\sqrt{\kappa_i}.
\]
Weights are $s^{o(1)}$ in these normalized units.
Equation~\eqref{eq:CE} and the terminal weight count force the hit
probability to be power-small.  The estimate is uniform in the
candidate slabs.  Taking measurable near maximizers, or a fine
finite net of their bounded coefficients with a slightly enlarged
width, proves Equation~\eqref{eq:rt-plane-nonconcentration}.

For waves take $v>0$ in the interior and put $j=1-2v$, so
$b<j<1$.  Ignore future planned insertions after $b$ and insert
one peel at $j$ on its free preserving continuation.  Its raw
assignment family saturates the segment bound and has
$r_j/R_v=s^{o(1)}$, $f_j=s^{o(1)}$ and the weight count of
Lemma~\ref{lem:rt-laws}.  Keep its assignment multipliers fixed
when comparing synthesized contributions.

For each $(C,C')$, greedily remove disjoint portions of the raw
$b$-rows contained in a slab and having mass greater than
$\rho^\tau$ relative to the original conditional law.  There
are at most $\rho^{-\tau}$ portions for each pair.  Stop when
every slab has residual original mass at most $\rho^\tau$.
The plane choices use a countable dense set and an arbitrarily small
enlargement of their widths.  Take their joint Borel versions under
the full analyzed $(C,C',b\text{-row})$ energy law, as in
Lemma~\ref{cyl:borel-paths}, before any later path restriction.

Split removed portions before synthesis by their pair and their
plane.  A label at $j$ receives contributions from only
$s^{-o(1)}$ name pairs: its angular width is $R_vs^{o(1)}$,
and Equation~\eqref{eq:acc} backflows its positions to the two
ancestor grids with small lists.  Coherent recombination of the
removed contributions at that label therefore costs at most
$s^{-o(1)}\rho^{-O(\tau)}$.  Indeed, after freezing its assignment
operator, if $v_{\lambda,\nu}$ are the analyzed contributions and
$N_\lambda\le s^{-o(1)}\rho^{-\tau}$ is their number, then
\[
 \left\|\sum_\nu v_{\lambda,\nu}\right\|_2^3
 \le N_\lambda^2\sum_\nu\|v_{\lambda,\nu}\|_2^3.
\]
Nonconnecting portions are
negligible by complementary projections.

For one removed component, insert a regularity split at
$i=j-h'>b$, using the appropriate step from
Proposition~\ref{prop:rt-plane}.  Its input at $b$ has regularity
$s^{-o(1)}\kappa_b$ and energy charged to its raw row mass
$e_{\mathrm{comp}}$.  A level of raw mass $z_i$ and threshold
$\widetilde\kappa_i$ satisfies
\[
 z_i\sqrt{\widetilde\kappa_i}
 \le s^{(1/2-\Gamma)(i-b)-o(1)}
                         e_{\mathrm{comp}}\sqrt{\kappa_b}
\]
by the prefix cube upper bound.  The low-floor remainder obeys the
same estimate when the floor is chosen sufficiently small.
Its synthesis has regularity at most
$s^{-o(1)}\widetilde\kappa_i$, with summed energy charged to
$z_i$.
In $Z$-coordinates, synthesis at $b$ and at $i$ broadens angular
spectral support only by $O(d_{\ell_b}+d_{\ell_i})/R_v$.
Since $H/(\ell_iR_v^2)=\Delta s^{o(1)}$, this is within the
$s^{-o(1)}\sqrt\Delta$ thickening permitted in
Proposition~\ref{prop:rt-plane}.  Changing angular units multiplies
input regularity by $R_v^4$, canceling the ending weight.
Relative to the ordinary free bound on the step $i\to j$, the tiny-cap
branch gains $\Gamma h'$, while the transverse branch gains
$(\Gamma-4\eta)h'-C\lambda h$, for a fixed $C$ absorbing the
$P^C$ factors. Both are positive by the choice of $h'$ and then
$\lambda$. The prefix $b\to i$ contributes only the baseline factor
$s^{(1/2-\Gamma)(i-b)-o(1)}$, so each branch retains its strict gain
over the full $b\to j$ comparison.
Sum the components, whose raw budgets sum to at most
$s^{-o(1)}E_b$, and the logarithmic levels.  Choose $\tau$ after
the two strict gains so that the recombination factor cannot absorb
them.  The removed field is then power-insufficient for the fixed
$j$-certificate.

The residual field must still have the full exponent of this
certificate.  The ordinary cube upper bound applied to its
synthesis at $b$, with inherited regularity, forces its raw mass
to be $s^{o(1)}E_b$.  Under the normalized residual law, the
conditional surviving fraction in a name pair is typically at
least $\rho^{\tau/2}$: pairs with smaller surviving fraction
contribute at most $\rho^{\tau/2}$ of the original total mass,
which is negligible relative to the subpower residual budget.
Every conditional residual plane mass is therefore at most
$\rho^{\tau/2}$ on the typical pairs.  This proves
Equation~\eqref{eq:rt-plane-nonconcentration}.  Deterministic
profiles and Lemma~\ref{lem:rt-time} survive by their subpower
robustness.  The removal need only be performed for the one finite
tangent query under consideration.
\end{proof}

\subsection{Differentiation and stationarity of the angular profile}

\begin{lemma}[Diagonal differentiation of an interval function]
\label{lem:rt-interval-density}
Let $J(x,y)$ be continuous for $x\le y$ in a compact interval,
with $0\le J(x,y)\le M(y-x)$ and
$J(x,z)\ge J(x,y)+J(y,z)$.  Then there is a measurable
$m\in[0,M]$ such that
\[
 \lim_{r\downarrow0}\frac{J(x,x+r)}r=m(x)
 \quad\hbox{for almost every }x.
\]
Applied to the wave profiles, there is a jointly measurable version
$m(a,v)\in[0,2]$, nonincreasing in $a$ for fixed $v$, for which
\begin{equation}\label{eq:rt-angular-tangent}
 \frac{J(a,v,v+zh)}h\longrightarrow z m(a,v)
\end{equation}
at almost every interior $(a,v)$, for each fixed $z>0$.
\end{lemma}

\begin{proof}
Define
\[
 \mu([x,y])=\inf_{\mathcal P}
                  \sum_{[u,w]\in\mathcal P}J(u,w),
\]
where $\mathcal P$ runs over finite partitions of $[x,y]$.
Joining two partitions proves one direction of additivity.
For the other, insert the joining point into any partition; by
superadditivity this can only decrease its sum.  Thus $\mu$ is
additive and $0\le\mu([x,y])\le M(y-x)$.  Its cumulative
function is nondecreasing and $M$-Lipschitz, so
$\mu([x,y])=\int_x^ym(t)\,dt$ for some $0\le m\le M$.

Set $K(x,y)=J(x,y)-\mu([x,y])$.  It is nonnegative and
superadditive, and its partition-sum infimum on every interval is
zero.  Fix $\epsilon_0>0$ and consider
$S=\{x:\limsup_{r\downarrow0}K(x,x+r)/r>\epsilon_0\}$.
For any $\delta_0>0$, choose a finite partition whose sum of
$K$ is less than $\delta_0$.  Away from its finitely many
endpoints, the intervals $[x,x+r]$ lying in the same partition cell
and satisfying $K(x,x+r)>\epsilon_0r$ form a Vitali cover of
$S$.  A disjoint countable subfamily covers $S$ up to a null set.
Superadditivity and nonnegativity, first for finite subfamilies
inside each cell, bound the sum of their $K$ values by
$\delta_0$.  Hence $|S|\le\delta_0/\epsilon_0$.
Let $\delta_0\downarrow0$, then use a countable sequence of
$\epsilon_0\downarrow0$.  The excess $K(x,x+r)/r$ tends to
zero almost everywhere.  Lebesgue differentiation of $\mu$
proves the first assertion with the full $r\downarrow0$ limit.

For the parameterized profiles define, on their common domain,
\[
 m(a,v)=\limsup_{\substack{r\downarrow0\\r\in\mathbb Q}}
                         \frac{J(a,v,v+r)}r.
\]
Continuity makes this jointly measurable, and monotonicity of $J$
in $a$ makes this version nonincreasing in $a$ pointwise.
The first assertion and Fubini identify it with the diagonal
derivative almost everywhere.  Substituting $r=zh$ gives
Equation~\eqref{eq:rt-angular-tangent}.
\end{proof}

\begin{lemma}[Dyadic stationarity]\label{lem:rt-stationarity}
At almost every interior wave point, and almost every classical
depth $a$ with $v=0$, along $h=2^{-n}\downarrow0$ one has
\begin{equation}
 J(a,v,v+Kh)-J(a+Kh,v,v+Kh)=o(h).
 \tag{stat}\label{eq:stat}
\end{equation}
\end{lemma}

\begin{proof}
Work on a compact rectangle inside the validity domain and let
$\delta=Kh$.  For fixed $v$, the function
$f(a)=J(a,v,v+\delta)$ is nonincreasing and lies between zero
and $2\delta$.  Consequently
\[
 \int_{a_0}^{a_1}\bigl(f(a)-f(a+\delta)\bigr)\,da
 =\int_{a_0}^{a_0+\delta}f(a)\,da
  -\int_{a_1}^{a_1+\delta}f(a)\,da
 \le2\delta^2.
\]
After integration in $v$, the integral of the nonnegative
left side of Equation~\eqref{eq:stat}, divided by $h$, is
$O(h)$.  Its sum over dyadic $h$ is finite, so Tonelli's theorem
implies convergence to zero almost everywhere.  A countable
exhaustion by compact rectangles proves the wave statement.
For classical $v=0$ omit the $v$ integral.
\end{proof}

Choose a point where the applicable conclusions of
Lemmas~\ref{lem:rt-interval-density} and
\ref{lem:rt-stationarity} hold and $D$ is differentiable.
Write $o_T=\partial_aD(a,v)$.  Classically use instead the constant
$m$ of Equation~\eqref{eq:rt-classical-m}.
For each sufficiently small dyadic $h$, choose the finite
configuration of Equation~\eqref{eq:query}, and perform
Lemma~\ref{lem:rt-slab-removal} before sampling additional rows.
All profile approximations at this query are taken before letting
$h$ tend to zero.  Rational perturbations of size $o(h)$ are
permitted by the Lipschitz profile bounds.

There is a bounded representative $O\in\R^4$, determined by
$(C,C')$, such that
\begin{equation}\label{eq:rt-incidence-point}
 O=(0,W)+\theta e(Z)+O(\rho^K),\qquad
 e(Z)=(1,Z,|Z|^2).
\end{equation}
Take the spacetime position at the finer ancestor and its time,
expressed in the moving frame of $C$.  Its discrepancy from the
later row position is at most $O(\rho^K)$, since both the time
discrepancy and the centered velocity are bounded at that scale.
Given $C$, the unit-depth name $O_1$ and $C(a+h,v)$ determine
one another to bounded lists: move between their times by at most
$O(\rho)$ with bounded centered velocity.  Thus
\begin{equation}\label{eq:rt-point-score}
 i_\rho(O_1\mid C)\longrightarrow o_T,
 \qquad
 \frac{H(O_1\mid C)}{\log(1/\rho)}\longrightarrow o_T.
\end{equation}
Here $H(\cdot\mid\cdot)$ denotes Shannon conditional entropy.
The angular statements give, for every needed finite $0<z\le1$,
\begin{equation}\label{eq:rt-angular-score}
 i_\rho(Z_z\mid C)\ge mz-o(1).
\end{equation}
For waves the expected score has limit $mz$.  Moreover
\begin{equation}\label{eq:rt-stationary-information}
 \frac{I(Z_K;C'\mid C)}{\log(1/\rho)}=o(1).
\end{equation}
Indeed $C'$ determines $C$ to bounded ambiguity, $Z_K$ is
equivalent given the roots to $A_{v+Kh}$, and the two conditional
angular entropies differ by the left side of
Equation~\eqref{eq:stat}, in $s$-units.  Both profiles remain
valid on the residual law by Lemma~\ref{lem:rt-stability}.
Equation~\eqref{eq:LP}, including the small-information version,
therefore preserves Equation~\eqref{eq:rt-angular-score} after
conditioning on $C'$.

\begin{lemma}[Width derivative]\label{lem:rt-width-derivative}
At almost every interior wave point,
\begin{equation}
 \partial_vD\le m+o_T-d+1.
 \tag{rv}\label{eq:rv}
\end{equation}
\end{lemma}

\begin{proof}
Given $C$, the bins $(Z,W)_1$ determine $C(a,v+h)$ except
for at most $O(\rho^{-1})$ possibilities.  The extra factor is
exactly the further vertical precision: the new vertical grid is
$\rho^2$ in the old units, while $W_1$ has width $\rho$.
The new angular center and the known horizontal bin determine its
change of tilt at the required accuracy.

On the other hand, adjoining $\theta_1$ to $(Z,W)_1$ adds
conditional entropy at least $(d-o(1))\log(1/\rho)$.
This follows from Lemma~\ref{lem:rt-time}, since $Y_\#$ is a
finer name than $(Z,W)_1$, and from entropy monotonicity.
Equation~\eqref{eq:rt-incidence-point} shows that
$(Z,W,\theta)_1$ is determined to bounded lists by $(Z_1,O_1)$
given $C$.  Its entropy is therefore at most
$(m+o_T+o(1))\log(1/\rho)$.
Subtract the time contribution, add the one vertical branching
exponent, and use the tangent increment of $D(a,v+h)-D(a,v)$.
This proves Equation~\eqref{eq:rv}.
\end{proof}

\subsection{The outgoing entropy inequality}

\begin{proposition}[Outgoing entropy]\label{prop:rt-outgoing}
Use the exponent data $d,p,\Gamma$ of Equation~\eqref{eq:rt-pd}.
For almost every classical depth $a\in(0,1)$ at which
$D(a,0)$ is differentiable and Lemma~\ref{lem:rt-stationarity}
applies, put $o_T=\partial_aD(a,0)$ and $m=3-p$.
For waves, let $(a,v)$ be almost every interior point with
$0<a<1$ and $0<v<(1-a)/2$ at which $D$ is differentiable and
Lemmas~\ref{lem:rt-interval-density} and \ref{lem:rt-stationarity}
apply, and put $o_T=\partial_aD(a,v)$ and $m=m(a,v)$.
In either case,
\begin{equation}
 o_T\ge4d+\min(1-d,m/2).
 \tag{round-out}\label{eq:round-out}
\end{equation}
\end{proposition}

\begin{proof}
Fix the finite incidence queries required by
Equations~\eqref{eq:inc}--\eqref{eq:cond} for a prescribed strict
final loss.  Later we choose the geometric and score tolerances
small compared with these queries.  Given $(C,C')$, sample $M$
independent copies of $(Z,W,\theta)$, all sharing the
representative $O$.  Let $L_i=(Z_i)_K$ be their fine angular
names, represented by bounded grid points, and put $e_i=e(L_i)$.
The integer $M$ will be fixed before taking the tangent limit.

The time-spread estimate survives conditioning on all these pins.
The own pin is already part of $Y_{\#,i}$, and conditional
independence gives
\begin{align*}
 I(\mathbf L_{\ne i};Y_{\#,i},\theta_i\mid C)
 &\le I(\mathbf L_{\ne i};C'\mid C)\\
 &\le\sum_{j\ne i}I(L_j;C'\mid C).
\end{align*}
For clarity, the second inequality follows by writing the mutual
information as conditional entropy minus entropy given $C,C'$:
the former is at most the sum of the marginal entropies, and the
latter equals their sum by conditional independence.
Equation~\eqref{eq:rt-stationary-information} makes the whole
information cost $o(\log(1/\rho))$ for fixed $M$.
Thus Lemma~\ref{lem:rt-time} and Equation~\eqref{eq:LP} give
the required time scores even with these pins.  The angular score
also satisfies, typically,
\begin{equation}\label{eq:rt-product-angular}
 i_\rho((Z_k)_z\mid C,C',\mathbf L_{\ne k})\ge mz-o(1),
\end{equation}
because the conditional law of copy $k$ is unchanged by the other
copies when $C,C'$ are given.

Use the law already obtained from
Lemma~\ref{lem:rt-slab-removal}, with a sufficiently small
$\eps>0$.  With probability tending to one, every four distinct
$e_i$ have determinant of magnitude at least $\rho^{C\eps}$,
and every pair of angular representatives has separation at least
$\rho^{C\eps}$.  To check this, sample successively in the
first-coordinate-one affine slice of $\R^4$.  The affine span
of the preceding points is contained in an affine plane in its
three remaining coordinates.  The probability that the next
point is within $\rho^\eps$ of such a plane tends to zero,
uniformly in the previously chosen points.  Successive distances
and bounded vector lengths give the determinant bound.  A plane
through one preceding point also bounds the probability of a
close pair; $V$ is Lipschitz on the bounded angular range.
Pin rounding by $O(\rho^K)$ is negligible.  There are finitely
many tuples, so all the conditions hold simultaneously.  We call
this the transverse event.

Four transverse pins already give the lower bound $4d$, for every
$m\ge0$.  Choose any four of the sampled pins and apply Gram--Schmidt
in their order to produce an orthonormal basis $(f_1,f_2,f_3,f_4)$,
renumbering these pins from $1$ to $4$ within this calculation.
Reveal the bins of $O\cdot f_i$ in reverse order.  At step $i$,
adjoin $Y_{\#,i}$ to the conditioning data.  It predicts all the
previously read components to bounded ambiguity, since their
directions are perpendicular to $e_i$ and
Equation~\eqref{eq:rt-incidence-point} eliminates their time
dependence.  The new component determines $\theta_i$ to at most
$\rho^{-C\eps}$ cells: its coefficient
$e_i\cdot f_i$ is bounded below by $\rho^{C\eps}$.
The time score with all pins consequently contributes at least
$d-C\eps-o(1)$ at each step.  The chain rule, the list bounds,
and Equation~\eqref{eq:LP} show that the score of $O_1$ is at
least $4d-C\eps-o(1)$.  The change from the orthonormal bins
to ordinary $O_1$ bins has bounded multiplicity given the pins.
Drop the pins by Equation~\eqref{eq:LP} and use
Equation~\eqref{eq:rt-point-score}; then let $\eps\downarrow0$.
This proves Equation~\eqref{eq:round-out} when $m=0$.

When $m>0$, we sharpen two of these four time contributions by
projecting onto a plane.  For $i<j$ let $\pi_{ij}$ be orthogonal
projection onto the two-dimensional quotient by
$\operatorname{span}(e_i,e_j)$, with measurable orthonormal
coordinates, and define the projective direction
\[
 g_{ij}(Z)=[\pi_{ij}e(Z)].
\]
Only values on the transverse event matter, so the definition at
zero may be arbitrary.  Color each triple $i<j<k$ by the vector
of quantized scores
\[
 i_\rho((g_{ij}(Z_k))_z\mid C,C',\mathbf L_{\ne k})
\]
at the finitely many required widths, with color intervals of
length $\eps$.  Use one overflow color.  It is avoided with
probability tending to one, since a bounded projective chart has
at most $O(\rho^{-z})$ cells and counting controls the upper
tail of its surprise.  There are a fixed finite number of colors.
The same finite Ramsey theorem \cite[Theorem~B, p.~267]{Ramsey1930}
therefore supplies, for sufficiently large fixed $M$, four indices whose
triples all have the same color vector.

Here is the geometric score constraint on such a quadruple.
Let $i<j<l<k$.  Given the two directions
$g_{ij}(Z_k)$ and $g_{il}(Z_k)$ to width $\rho^z$, and all
pins except $L_k$, the number of possible good $(Z_k)_z$ cells
is at most $\rho^{-C\eps}$.  Fix one compatible transverse
witness $Z_0$.  The first directional constraint places $e(Z)$
within $O(\rho^{z-C\eps})$ of
$\operatorname{span}(e_i,e_j,e(Z_0))$, and the second does the
same for $\operatorname{span}(e_i,e_l,e(Z_0))$.
The determinant lower bound controls the inverse linear algebra,
so their intersection confines $e(Z)$ within
$O(\rho^{z-C\eps})$ of
$\operatorname{span}(e_i,e(Z_0))$.
In the first-coordinate-one slice this is the secant line
\[
 (1-t)e(L_i)+t e(Z_0).
\]
Its quadratic null residual is
$t(1-t)|L_i-Z_0|^2$ (up to the harmless choice of sign).
The roots $0,1$ are separated and their derivatives have magnitude
at least $\rho^{C\eps}$.  Boundedness and the pair separation
therefore leave only neighborhoods of radius
$O(\rho^{z-C\eps})$ of the two endpoints.  Each contains at
most $\rho^{-C\eps}$ angular $\rho^z$-cells.
This proves the list bound, also when $z\le C\eps$, by the
trivial ambient angular count.

Apply this list bound and Equation~\eqref{eq:rt-product-angular}
with the same conditioning data $C,C',\mathbf L_{\ne k}$.
The joint score of the two directions is at least
$mz-C\eps-o(1)$.  Their marginal score sum is at least their
joint score up to $o(1)$ outside a vanishing exceptional set,
by Equation~\eqref{eq:LP} and the chain rule.  Thus
\begin{equation}\label{eq:rt-two-directions}
 i_\rho((g_{ij}(Z_k))_z\mid C,C',\mathbf L_{\ne k})
 +i_\rho((g_{il}(Z_k))_z\mid C,C',\mathbf L_{\ne k})
 \ge mz-C\eps-o(1).
\end{equation}
On a monochromatic quadruple these two scores differ by at most
$\eps$.  Hence both are at least
$(m/2)z-C\eps-o(1)$ simultaneously at every required width.
There are finitely many possible quadruples; after passing to a
subsequence, select fixed indices $i<j<k$ and an event of fixed
positive probability on which all these
bounds hold.  This restriction has probability bounded below
independently of $\rho$ after the parameters are fixed, so the
finite score bounds and the small-information consequences survive
by likelihood comparison.  Retain transversality as well.

Use the side data $D_*=(C,\mathbf L_{\ne k})$, the point
$Q=\pi_{ij}O$, and the line $\mathcal L$ through
$\pi_{ij}(0,W_k)$ in direction $\pi_{ij}e(Z_k)$.
They are incident up to $O(\rho^K)$.  Restrict to a
positive-probability bounded slope chart after a fixed rotation;
a finite collection of such charts covers all directions.  We
verify Equation~\eqref{eq:cond} with exponents $d$ and $m/2$.

For the point score, adjoin $Y_{\#,k}$.  This determines the
line at the required conditioning resolution to at most
$\rho^{-C\eps}$ choices, by transversality and fine pinning.
Conditioning additionally on its line cell therefore costs at
most $C\eps+o(1)$ in score.  Given these data, the projected
point determines $\theta_k$ at width $\rho^{z'}$ to at most
$\rho^{-C\eps}$ possibilities, because the projected speed is
at least $\rho^{C\eps}$.  Lemma~\ref{lem:rt-time}, its
preservation under the other pins, and Equation~\eqref{eq:LP}
give
\[
 i_\rho(Q_{z'}\mid\mathcal L_z,D_*)
                 \ge dz'-C\eps-o(1),\qquad z'\le z.
\]
For the line score, adjoin $C'$.  The point is then known, and
the line cell determines its direction at the same resolution,
up to the chart's controlled lists.  The directional lower bound
from Equation~\eqref{eq:rt-two-directions}, with the same
conditioning $C,C',\mathbf L_{\ne k}$, gives
\[
 i_\rho(\mathcal L_{z'}\mid Q_z,D_*)
                 \ge (m/2)z'-C\eps-o(1).
\]
These are exactly the finite-grid hypotheses of
Equation~\eqref{eq:cond}.  Choose $\eps$ and the earlier
tolerances sufficiently small.  Equation~\eqref{eq:inc} yields
\begin{equation}\label{eq:rt-projected-score}
 i_\rho(Q_1\mid D_*)
                 \ge d+\min(1,d+m/2)-o(1)
\end{equation}
to any prescribed strict loss.

It remains to recover the other two coordinates of $O$.  Let
$f_i$ be the unit vector along $e_i$, and let $f_j$ be the unit
vector along the component of $e_j$ perpendicular to $e_i$.
After $Q_1$, reveal $(O\cdot f_j)_1$, then
$(O\cdot f_i)_1$.  For the former, adjoin $Y_{\#,j}$; it
predicts $Q_1$ to bounded ambiguity because $\pi_{ij}e_j=0$,
and the new coordinate reveals $\theta_j$ to
$\rho^{-C\eps}$ choices.  For the latter, adjoin
$Y_{\#,i}$; it predicts both preceding coordinates because
they annihilate $e_i$, and this last coordinate reveals
$\theta_i$.  Each conditional score is therefore at least
$d-C\eps-o(1)$ by the time estimate with the pin data.
All these conditioning variables are available: $D_*$ includes
the own pins $L_i,L_j$, and the information argument at the
start of the proof permits the other pins.

Combining these two contributions with
Equation~\eqref{eq:rt-projected-score} and then dropping $D_*$
down to $C$ by Equation~\eqref{eq:LP} gives
\[
 i_\rho(O_1\mid C)
 \ge3d+\min(1,d+m/2)-o(1)
 =4d+\min(1-d,m/2)-o(1)
\]
on the retained event of fixed positive probability.
Its deterministic limit in Equation~\eqref{eq:rt-point-score}
forces the asserted inequality.

The order of choices is essential.  Fix the desired final strict
loss and the finite grids in Equation~\eqref{eq:inc}; next choose
the geometric and color tolerances; next fix a Ramsey-sufficient
$M$; then let dyadic $h\downarrow0$ through the typical point,
taking all previous finite-list and inner limits sufficiently
accurate for that $h$.  Every restriction made after the
stationarity estimate has a positive probability bound at these
fixed choices.  No arbitrary subpower restriction is used to
preserve a mutual-information assertion.
\end{proof}

\subsection{The differential contradiction}

\begin{lemma}[A two-direction descent]\label{lem:rt-descent}
Let $\Omega=\{(t,v):0<t<1,\ 0<v<t/2\}$.  Suppose that
$E$ is locally Lipschitz, nonnegative, and has zero trace on
$v=t/2$.  Let $m$ be measurable and nondecreasing in $t$ at
each fixed $v$, with this condition allowed outside a null set of
$v$.  If $c>0$ and $q\in\R$, the following pair of almost
everywhere inequalities is impossible:
\[
 E_t+\tfrac12E_v\le-c\quad\hbox{where }m\le q,
 \qquad
 E_t\le-c\quad\hbox{where }m\ge q.
\]
\end{lemma}

\begin{proof}
Fix $0<t_0<t_1<1$, put $T=t_1-t_0$, and let $N$ be the
null set where differentiation or an applicable inequality fails.
Fubini after the affine change
$(t,\delta)\mapsto(t,t/2-\delta)$ shows that almost every
line $\gamma_\delta(t)=(t,t/2-\delta)$ meets $N$ in a null
set.  Ordinary Fubini gives the same property for almost every
vertical section.  Include exceptional monotonicity sections in
the latter null set.  Choose a good $\delta>0$ sufficiently
small that $\delta<t_0/4$ and
$E_0:=E(t_0,t_0/2-\delta)<cT/2$; the trace condition permits
this.  The function $f(t)=E(\gamma_\delta(t))$ is Lipschitz
on $[t_0,t_1]$; let $B$ be a Lipschitz constant.

Set $A=\{t:m(\gamma_\delta(t))>q\}$.  If $|A|=0$, the
first inequality integrates to $f(t_1)\le E_0-cT<0$.
Otherwise let $\tau$ be the essential first time in $A$.
Then $\tau<t_1$, $A$ has zero measure before $\tau$, and
every sufficiently short interval immediately after $\tau$
has positive intersection with $A$.  Integration before $\tau$
gives $f(\tau)\le E_0-c(\tau-t_0)$.
Choose $t_*\in A$ arbitrarily close to $\tau$ from above,
with $v_*=t_*/2-\delta$ a good vertical section.  Such a
choice exists because this affine map of $t_*$ preserves null
sets.  Monotonicity gives $m(t,v_*)>q$ for $t\ge t_*$,
so the second inequality integrates along this vertical segment:
\[
 E(t_1,v_*)\le f(t_*)-c(t_1-t_*)
 \le E_0-cT+(B+c)(t_*-\tau).
\]
Taking $t_*$ sufficiently close to $\tau$ makes the right side
negative, again a contradiction.  Each segment is a compact subset
of $\Omega$, so local Lipschitz regularity suffices.
\end{proof}

\begin{proof}[Proof of Theorem~\ref{thm:round}]
In the classical contradiction,
Equation~\eqref{eq:rt-classical-m} gives
$m/2=1-d+\Gamma$.  Proposition~\ref{prop:rt-outgoing}
therefore implies
$\partial_aD(a,0)\ge1+3d$ almost everywhere.
The profile is Lipschitz and $D(0,0)\ge0$ by continuity of the
lower bound in Proposition~\ref{prop:rt-profiles}.  Integration
gives $D(1,0)\ge1+3d$.  But the endpoint bound gives
\[
 D(1,0)\le p+d=1+3d-2\Gamma,
\]
a contradiction.  Hence
$\Gamma_{\mathrm{cl},\mathrm{round}}\le100\eta$.

For waves suppose the exponent is larger than $100\eta$; the
classical conclusion just obtained supplies the strict comparison
needed in Proposition~\ref{prop:rt-planck}.  Put $t=1-a$ and
use the same letter $E$ for the excess in these coordinates.
It is Lipschitz, nonnegative, and zero on $v=t/2$.
The version of $m$ from Lemma~\ref{lem:rt-interval-density} is
nondecreasing in $t$ at fixed $v$.  Equations~\eqref{eq:rv}
and \eqref{eq:round-out} imply
\[
 E_t=1+3d-2\Gamma-o_T,
 \qquad E_v\le m+o_T-d-3.
\]
If $m\le2-2d$, then
$o_T\ge4d+m/2\ge m+5d-1$, so
\[
 E_t+\tfrac12E_v\le-2\Gamma.
\]
If $m\ge2-2d$, then $o_T\ge1+3d$, so
$E_t\le-2\Gamma$.
Apply Lemma~\ref{lem:rt-descent} with
$c=2\Gamma>0$ and $q=2-2d$.  This is a contradiction.
The free wave exponent is therefore at most $100\eta$.
The finite-stack comparison completes all assertions of the theorem.
\end{proof}
